%%%%%%%%%%%%
%
% $Autor: Wings $
% $Datum: 2019-03-05 08:03:15Z $
% $Pfad: Functions.tex $
% $Version: 4250 $
% !TeX spellcheck = en_GB/de_DE
% !TeX encoding = utf8
% !TeX root = manual 
% !TeX TXS-program:bibliography = txs:///biber
%
%%%%%%%%%%%%

\chapter{Funktionen}

In diesem Kapitel werden die einzelnen Funktionen von ELMAR beschrieben.

\section{Referenzfahrt/Homing-Funktion}

Die Referenzfahrt, dient zur Initialisierung der Achsen von ELMAR. Dabei fahren die Achsen kontrolliert ihre Referenzpunkte an den Endschaltern an. Nach erfolgreicher Referenzierung kennt das System die Position aller Achsen im Maschinenkoordinatensystem.

Anschließend fährt das System auf die vorprogrammierte Home-Position. Diese befindet sich 7 mm von den Endschaltern entfernt.

\section{Start-Funktion}

Durch einfaches Drücken des Start-Knopfes wird der Schneidvorgang gestartet oder nach Pausieren fortgesetzt. 
Der Heizdraht wird aktiviert und die Achsen beginnen mit der Ausführung der geladenen Schneiddatei.

\section{Pause-Funktion}

Mit dem Pause-Knopf kann der Schneidvorgang jederzeit unterbrochen werden. 
Die Bewegung der Achsen wird gestoppt, während der aktuelle Zustand beibehalten wird.

\section{Stopp-Funktion}

Das Drücken des Stopp-Knopfes quittiert den Schneidauftrag vollständig.

\section{Homing-Funktion}
 
Die Homing-Funktion kann in UGS durch die Eingabe des Befehls \$H ausgeführt werden.
\section{Not-Aus-Funktion}

Der Not-Aus-Schalter dient der sofortigen Unterbrechung aller Funktionen. 
Im Notfall werden sämtliche Bewegungen gestoppt und die Stromversorgung wird unterbrochen.